%%%%%%%%%%  scan idx 30  |  folio 17  %%%%%%%%%%
tout comme la cohomologie et l'homologie. Ainsi, si $f:X\to Y$ est un morphisme de sch\'emas alg\'ebriques propres, les deux homomorphismes
\[
f^*:\Gr^\bullet(Y)\longrightarrow \Gr^\bullet(X)\qquad\text{et}\qquad f_*:\Gr_\bullet(X)\longrightarrow \Gr_\bullet(Y)
\]
sont formellement transpos\'es l'un de l'autre, au sens des accouplements (6.6).

Les homomorphismes (6.4) et (6.5) et l'accouplement (6.6) qu'on en d\'eduit, peuvent \^etre consid\'er\'es comme la source des \uline{relations num\'eriques} en th\'eorie des intersections sur les sch\'emas alg\'ebriques propres sur un corps donn\'e $k$. Ce point de vue sera d\'evelopp\'e avec quelques d\'etails dans Exp.\ X, ind\'ependant pour l'essentiel des Exp.\ VI, VII, VIII (centr\'es sur la d\'emonstration du th\'eor\`eme de Riemann--Roch) et de Exp.\ IX. On y donnera \'egalement quelques notions sur le \uline{comportement de} $K^\bullet(X),\Gr^\bullet(X)$ \uline{par sp\'ecialisation}.

\subsection*{7.}
Comme dans toute th\'eorie d'intersection, il importe de pr\'eciser la place que prend dans la th\'eorie le \uline{groupe de Picard}, (ou groupe des \uline{classes de diviseurs} dans la terminologie classique). Moyennant une l\'eg\`ere restriction, automatiquement satisfaite lorsque $X$ est par exemple quasi-compact, on \'etablit dans Exp.\ X un isomorphisme canonique
\begin{equation}\tag{7.1}
c_1:\Pic(X)\xrightarrow{\sim}\Gr^1(X),
\end{equation}
qui montre le r\^ole exact du groupe de Picard dans la th\'eorie que nous proposons. Ce r\^ole justifie une \'etude plus d\'etaill\'ee de ce groupe de Picard dans les expos\'es XII, XIII, qui sont logiquement ind\'ependants du texte du S\'eminaire. Une \'etude de l'\uline{\'equivalence num\'erique} et de certains \uline{th\'eor\`emes de finitude pour le groupe de Picard} (annonc\'es sans d\'emonstration dans [5]) est faite par Kleiman dans Exp.\ XIII, en utilisant des techniques de nature purement num\'erique (n'utilisant pas la th\'eorie du foncteur de Picard et de sa repr\'esentabilit\'e), dues \`a lui-m\^eme et \`a Mumford, nettement plus simples que les d\'emonstrations initiales de Grothendieck. Le seul r\'esultat r\'efractaire aux m\'ethodes de Kleiman (et dont il est oblig\'e de faire usage) est [5, C-07 (i)], dont la d\'emonstration et certains raffinements (incluant certains th\'eor\`emes de repr\'esentabilit\'e pour le foncteur de Picard) occupent l'expos\'e XII. Ce dernier rel\`eve d'ailleurs d'un esprit et d'une technique enti\`erement \'etrangers au reste du S\'eminaire, se rattachant \`a [5]~; cet ex%
